
We present evidence for the semantic invalidity hypothesis in three parts: (1) differential effect confirmation, (2) dose-response relationship, and (3) rotation control validation.

\subsubsection{Differential Effect: Asymmetric Degradation, Symmetric Stability}

At flip probability p=0.5, asymmetric digits {2,3,5,6,7,9} consistently degrade 0.72-1.00 percentage points relative to baseline, while symmetric digits {0,1,8} remain stable with changes <0.2\%. This pattern replicates across four independent experiments.

\textbf{Table 1: Differential Effect Confirmation Across Sub-Hypotheses}

\begin{table*}[htbp]
\centering
\resizebox{\dimexpr 2\columnwidth+\columnsep\relax}{!}{%
\begin{tabular}{llllllllll}
\toprule
Hypothesis & Type & Seeds & Baseline Asym & Flip50 Asym & Asym $\Delta$ & Baseline Sym & Flip50 Sym & Sym $\Delta$ & Differential \\
\midrule
h-e1 & EXISTENCE & n=1 & 98.95\% & 98.23\% & -0.72\% & 99.43\% & 99.38\% & -0.05\% & 0.67\% \\
h-m1 & MECHANISM & n=5 & 99.02 $\pm$ 0.23\% & 98.24 $\pm$ 0.05\% & -0.78\% & Stable & Stable & ~0\% & 0.78\% \\
h-m & EXTENDED & n=5 & 98.99 $\pm$ 0.04\% & 97.99 $\pm$ 0.12\% & -1.00\% & 99.50 $\pm$ 0.08\% & 99.34 $\pm$ 0.04\% & -0.16\% & 0.84\% \\
\bottomrule
\end{tabular}
}
\end{table*}

\textit{Note: Differential = |Asym $\Delta |$ - |Sym $\Delta |$. Asymmetric digits degrade 3-$15\times$ more than symmetric digits across all validations. At extreme flip rate p=0.9, symmetric digits show slight degradation (-0.28\% to -0.77\%), indicating a boundary condition where general augmentation effects emerge at very high flip probabilities.}

The consistency across experiments is striking: asymmetric degradation ranges 0.72-1.00\% (within 0.3 pp despite different implementations), while symmetric changes remain <0.2\% at flip50 in all cases. Multi-seed validation in h-m1 and h-m shows tight seed variability (standard deviation 0.04-0.23\%), indicating reproducibility.

\subsubsection{Dose-Response: Deterministic Label Noise Mechanism}

The dose-response relationship shows asymmetric digit accuracy decreases monotonically across flip probabilities {0.0, 0.3, 0.5, 0.9}.

\textbf{Table 2: Dose-Response Statistical Tests}

\begin{table*}[htbp]
\centering
\resizebox{\dimexpr 2\columnwidth+\columnsep\relax}{!}{%
\begin{tabular}{lllll}
\toprule
Hypothesis & Flip Probabilities & Spearman $\rho$ & p-value & Interpretation \\
\midrule
h-m1 & {0.0, 0.3, 0.5, 0.9} & **-1.0000** & **<0.001** & Perfect monotonicity \\
h-m & {0.0, 0.3, 0.5, 0.9} & **-0.969** & **1.$97\times 10^{-12}$** & Near-perfect monotonicity \\
\bottomrule
\end{tabular}
}
\end{table*}

\textit{Note: Perfect $\rho$=-1.0 is exceptionally rare in empirical studies, indicating deterministic relationship. Both tests use n=20 data points (4 dose levels $\times$ 5 seeds).}

h-m1 achieves $\rho$=-1.0000 (mathematically perfect negative correlation), while h-m achieves $\rho$=-0.969 (very strong negative correlation). Both have p-values orders of magnitude below the significance threshold.

The perfect correlation in h-m1 merits explanation. Spearman $\rho$=-1.0 means zero rank inversions: every increment in flip probability corresponds to a decrease in asymmetric accuracy across all 20 seed-level measurements. This is extraordinary in noisy empirical machine learning experiments. We attribute this to three factors: (1) deterministic label noise mechanism (noise proportion = flip\_probability $\times$ asymmetric\_fraction), (2) low MNIST noise floor (high baseline accuracy ~99\% indicates task is well within model capacity), and (3) multi-seed averaging (n=5 seeds per dose level smooths random variation, with observed seed standard deviation <0.12\% being much smaller than dose-level gaps of 0.5-3.0 pp).

\textbf{Degradation Magnitude Across Flip Probabilities}: At flip probability p=0.3, asymmetric accuracy degrades 0.37-0.51 pp (h-m1/h-m). At p=0.5, degradation increases to 0.72-1.00 pp. At extreme p=0.9, degradation reaches 3.15-4.10 pp.

\subsubsection{Rotation Control: Isolating Semantic Invalidity}

To confirm that semantic invalidity (not general augmentation) is the causal mechanism, we compare rotation $\pm 15$° augmentation against baseline and flip conditions. Rotation causes no differential effect on asymmetric versus symmetric digits.

\textbf{Table 3: Rotation Control Validation}

\begin{table*}[htbp]
\centering
\resizebox{\dimexpr 2\columnwidth+\columnsep\relax}{!}{%
\begin{tabular}{llllll}
\toprule
Hypothesis & Condition & Asym Acc & Sym Acc & Asym vs Baseline & Group Diff (Sym-Asym) \\
\midrule
h-e1 & Baseline & 98.95\% & 99.43\% & --- & -0.48\% \\
h-e1 & Rotation & 99.14\% & 99.43\% & +0.19\% & -0.29\% \\
h-c1 & Baseline & 98.96\% & 99.35\% & --- & -0.39\% \\
h-c1 & Rotation & 99.10\% & 99.63\% & +0.14\% & -0.53\% \\
h-m & Baseline & 98.99 $\pm$ 0.04\% & 99.50 $\pm$ 0.08\% & --- & -0.51\% \\
h-m & Rotation & 99.04 $\pm$ 0.05\% & 99.45 $\pm$ 0.05\% & +0.05\% & -0.41\% \\
\bottomrule
\end{tabular}
}
\end{table*}

\textit{Note: ''Asym vs Baseline'' shows rotation effect on asymmetric digits. All rotation effects are <1\%, and group differentials are comparable between baseline and rotation (<0.2\% difference).}

Across three independent validations (h-e1, h-c1, h-m), rotation augmentation changes asymmetric digit accuracy by -0.14\% to +0.19\%---all within the 1\% threshold for practical equivalence. The group differential (symmetric minus asymmetric accuracy) differs by only 0.05-0.19\% between baseline and rotation conditions, far below the 1\% threshold.

\textbf{Comparison: Flip versus Rotation}: At flip probability p=0.5, asymmetric digits degrade 0.72-1.00\%. In contrast, rotation causes asymmetric changes of -0.14\% to +0.19\%---a 3-$20\times$ difference in effect magnitude. This stark contrast confirms that semantic invalidity is the distinguishing factor.

\subsubsection{Per-Digit Heterogeneity}

Per-digit analysis from h-m reveals heterogeneity within the asymmetric digit group. At flip probability p=0.9, digit 7 shows minimal degradation (-0.30\%), while digits 2 and 5 show severe degradation (-6.60\% and -6.93\%, respectively). Symmetric digits {0,1,8} cluster near zero degradation (-0.28\% to -0.77\%).

This heterogeneity suggests semantic validity operates on a continuum. Degradation magnitude may correlate with ''semantic distance''---the perceptual dissimilarity between a flipped digit and its canonical form. Flipped digit 7 may visually resemble canonical 7, resulting in low label noise and minimal degradation. In contrast, flipped digits 2 and 5 look distinctly different from their canonical forms, creating high label noise and severe degradation.

\textbf{Symmetric Digit Stability: Boundary Conditions at Extreme Flip Rates}: At moderate flip rates (p $\leq$ 0.5), symmetric degradation is -0.05\% to -0.16\%, confirming <0.2\% stability threshold. At extreme flip rate p=0.9, symmetric degradation is -0.28\% to -0.77\%, exceeding the threshold. This indicates general augmentation effects (independent of semantic validity) become non-negligible at very high flip probabilities.

\subsubsection{Cross-Hypothesis Consistency}

The high consistency across four independent sub-hypotheses validates reproducibility:

\begin{itemize}
\item \textbf{Effect Size (Flip50 Asymmetric Degradation)}: h-e1: -0.72\%, h-m1: -0.78\%, h-m: -1.00\%. Range 0.72-1.00\% represents good agreement.
\item \textbf{Dose-Response}: h-m1 achieved perfect $\rho$=-1.0, h-m achieved $\rho$=-0.969. Both confirm perfect/near-perfect monotonicity.
\item \textbf{Rotation Control}: All experiments show asymmetric changes <1\% threshold.
\item \textbf{Symmetric Stability at Moderate Flip Rates}: All experiments show <0.2\% change at flip50.
\end{itemize}

